% GENERATED FILE. Edit canonical lessons, then run npm run generate:books.
% canonical-source-hash: fnv1a64:52ee489b9fb87387
% canonical-lessons: RU-C305-vospityvat, RU-C305-usynovlyat, RU-C305-staret, RU-C305-vypusknoy, RU-C305-pokoleniye

\chapter{Bringing up a family}
\label{ch:ru-bringing-up-a-family}
\hlchaptermodality{305}

\begin{chapteropening}
\textbf{By the end of this chapter:} \emph{I can say that someone brought up or adopted a child or is getting old, and talk about a graduation party and a generation.}

Complete the last lesson of chapter 305: I can say that someone brought up or adopted a child or is getting old, and talk about a graduation party and a generation.
\end{chapteropening}

\section[\texorpdfstring{\ru{воспитывать} (vospítyvat)}{воспитывать (vospítyvat)}]{\ru{воспитывать} (vospítyvat) — to bring up, to raise}

\label{lesson:RU-C305-vospityvat}



\begin{quote}
Before the new one: say the Russian for old age, then the Russian for to grow up.
\end{quote}

\subsection*{You'll want to know: \ru{воспитывать}}
\textbf{\ru{воспитывать}} (\emph{vospítyvat}) — \textquotedblleft{}to bring up, to raise\textquotedblright{}.

\begin{quote}
\ru{Бабушка} \ru{воспитала} \ru{троих} \ru{внуков}.
\end{quote}

\begin{quote}
\emph{Bábushka vospitála troíkh vnúkov.}
\end{quote}

\begin{quote}
Grandma brought up three grandchildren.
\end{quote}

Imperfective; perfective \emph{\ru{воспитать}} (\emph{vospitát}).

\begin{cousinweb}
\emph{\ru{воспитывать}} is usually described as a Church Slavonic formation, from \emph{\ru{вос}-} \textquotedblleft{}up\textquotedblright{} and \emph{\ru{питать}} \textquotedblleft{}to feed, to nourish\textquotedblright{}: to feed up. EtymDB lists no parent.
\end{cousinweb}

\begin{grammarlens}[title={What you've built}]
One more word for telling, explaining and describing at B1.
\end{grammarlens}

\subsection*{Your turn}
\begin{itemize}
\raggedright
  \item \emph{Say it:} \emph{vospítyvat}
  \item \emph{Say it:} \emph{vospítyvat}, once more
  \item \emph{Say it:} a sentence of your own with \emph{vospítyvat}
  \item \emph{Recall:} say the Russian for old age, then the Russian for to grow up, then say \emph{vospítyvat} again
\end{itemize}

\begin{tcolorbox}[breakable,colback=teal!4,colframe=teal!35!black,title={Before you move on}]
What does \ru{воспитывать} mean? (\textquotedblleft{}To bring up, to raise\textquotedblright{}.) Say it once more.
\end{tcolorbox}
\section[\texorpdfstring{\ru{усыновлять} (usynovlyát)}{усыновлять (usynovlyát)}]{\ru{усыновлять} (usynovlyát) — to adopt (a child)}

\label{lesson:RU-C305-usynovlyat}



\begin{quote}
Before the new one: say the Russian for to grow up, then the Russian for to bring up.
\end{quote}

\subsection*{You'll want to know: \ru{усыновлять}}
\textbf{\ru{усыновлять}} (\emph{usynovlyát}) — \textquotedblleft{}to adopt (a child)\textquotedblright{}.

\begin{quote}
\ru{Они} \ru{усыновили} \ru{мальчика} \ru{из} \ru{детского} \ru{дома}.
\end{quote}

\begin{quote}
\emph{Oní usynovíli málchika iz détskovo dóma.}
\end{quote}

\begin{quote}
They adopted a boy from a children's home.
\end{quote}

Imperfective; the perfective is \emph{\ru{усыновить}} (\emph{usynovít}). For a girl also \emph{\ru{удочерить}}.

\begin{cousinweb}
\emph{\ru{усыновлять}} is built on \emph{\ru{у}-} and \emph{\ru{сын}} \textquotedblleft{}son\textquotedblright{}, inherited from Proto-Slavic \emph{synŭ}, from PIE \emph{suHnús} (EtymDB): to make someone a son.
\end{cousinweb}

\begin{grammarlens}[title={What you've built}]
One more word for telling, explaining and describing at B1.
\end{grammarlens}

\subsection*{Your turn}
\begin{itemize}
\raggedright
  \item \emph{Say it:} \emph{usynovlyát}
  \item \emph{Say it:} \emph{usynovlyát}, once more
  \item \emph{Say it:} a sentence of your own with \emph{usynovlyát}
  \item \emph{Recall:} say the Russian for to grow up, then the Russian for to bring up, then say \emph{usynovlyát} again
\end{itemize}

\begin{tcolorbox}[breakable,colback=teal!4,colframe=teal!35!black,title={Before you move on}]
What does \ru{усыновлять} mean? (\textquotedblleft{}To adopt (a child)\textquotedblright{}.) Say it once more.
\end{tcolorbox}
\section[\texorpdfstring{\ru{стареть} (starét)}{стареть (starét)}]{\ru{стареть} (starét) — to grow old, to age}

\label{lesson:RU-C305-staret}



\begin{quote}
Before the new one: say the Russian for to bring up, then the Russian for to adopt.
\end{quote}

\subsection*{You'll want to know: \ru{стареть}}
\textbf{\ru{стареть}} (\emph{starét}) — \textquotedblleft{}to grow old, to age\textquotedblright{}.

\begin{quote}
\ru{Мой} \ru{пёс} \ru{стареет} \ru{и} \ru{много} \ru{спит}.
\end{quote}

\begin{quote}
\emph{Moy pyos staréyet i mnógo spit.}
\end{quote}

\begin{quote}
My dog is getting old and sleeps a lot.
\end{quote}

Imperfective; perfective \emph{\ru{постареть}} (\emph{postarét}). Of things: \emph{\ru{устареть}} — to go out of date.

\begin{cousinweb}
\emph{\ru{стареть}} is built on \emph{\ru{старый}} \textquotedblleft{}old\textquotedblright{}, inherited from Proto-Slavic \emph{starŭ} \textquotedblleft{}old\textquotedblright{} (EtymDB), with the verb ending \emph{-\ru{еть}} for becoming.
\end{cousinweb}

\begin{grammarlens}[title={What you've built}]
One more word for telling, explaining and describing at B1.
\end{grammarlens}

\subsection*{Your turn}
\begin{itemize}
\raggedright
  \item \emph{Say it:} \emph{starét}
  \item \emph{Say it:} \emph{starét}, once more
  \item \emph{Say it:} a sentence of your own with \emph{starét}
  \item \emph{Recall:} say the Russian for to bring up, then the Russian for to adopt, then say \emph{starét} again
\end{itemize}

\begin{tcolorbox}[breakable,colback=teal!4,colframe=teal!35!black,title={Before you move on}]
What does \ru{стареть} mean? (\textquotedblleft{}To grow old, to age\textquotedblright{}.) Say it once more.
\end{tcolorbox}
\section[\texorpdfstring{\ru{выпускной} (vypusknóy)}{выпускной (vypusknóy)}]{\ru{выпускной} (vypusknóy) — a graduation party, a prom}

\label{lesson:RU-C305-vypusknoy}



\begin{quote}
Before the new one: say the Russian for to adopt, then the Russian for to grow old.
\end{quote}

\subsection*{You'll want to know: \ru{выпускной}}
\textbf{\ru{выпускной}} (\emph{vypusknóy}) — \textquotedblleft{}a graduation party, a prom\textquotedblright{}.

\begin{quote}
\ru{На} \ru{выпускном} \ru{все} \ru{танцевали} \ru{до} \ru{утра}.
\end{quote}

\begin{quote}
\emph{Na vypusknóm vse tantseváli do utrá.}
\end{quote}

\begin{quote}
At the graduation party everyone danced until morning.
\end{quote}

Short for \emph{\ru{выпускной} \ru{вечер}}; an adjective that serves as a noun.

\begin{cousinweb}
\emph{\ru{выпускной}} is built on \emph{\ru{выпуск}} \textquotedblleft{}a graduating class\textquotedblright{}, from \emph{\ru{выпускать}} \textquotedblleft{}to let out\textquotedblright{}, whose root is Proto-Slavic \emph{pustiti} \textquotedblleft{}to let go\textquotedblright{} (EtymDB).
\end{cousinweb}

\begin{grammarlens}[title={What you've built}]
One more word for telling, explaining and describing at B1.
\end{grammarlens}

\subsection*{Your turn}
\begin{itemize}
\raggedright
  \item \emph{Say it:} \emph{vypusknóy}
  \item \emph{Say it:} \emph{vypusknóy}, once more
  \item \emph{Say it:} a sentence of your own with \emph{vypusknóy}
  \item \emph{Recall:} say the Russian for to adopt, then the Russian for to grow old, then say \emph{vypusknóy} again
\end{itemize}

\begin{tcolorbox}[breakable,colback=teal!4,colframe=teal!35!black,title={Before you move on}]
What does \ru{выпускной} mean? (\textquotedblleft{}A graduation party, a prom\textquotedblright{}.) Say it once more.
\end{tcolorbox}
\section[\texorpdfstring{\ru{поколение} (pokoléniye)}{поколение (pokoléniye)}]{\ru{поколение} (pokoléniye) — a generation}

\label{lesson:RU-C305-pokoleniye}



\begin{quote}
Before the new one: say the Russian for to grow old, then the Russian for a graduation party.
\end{quote}

\subsection*{You'll want to know: \ru{поколение}}
\textbf{\ru{поколение}} (\emph{pokoléniye}) — \textquotedblleft{}a generation\textquotedblright{}.

\begin{quote}
\ru{Три} \ru{поколения} \ru{семьи} \ru{жили} \ru{в} \ru{одном} \ru{доме}.
\end{quote}

\begin{quote}
\emph{Tri pokoléniya semí zhíli v odnóm dóme.}
\end{quote}

\begin{quote}
Three generations of the family lived in one house.
\end{quote}

Neuter.

\begin{cousinweb}
\emph{\ru{поколение}} is built on \emph{\ru{колено}} \textquotedblleft{}knee\textquotedblright{}, which in older Russian also meant \textquotedblleft{}a branch of a family, a line of descent\textquotedblright{}; EtymDB lists no parent.
\end{cousinweb}

\begin{grammarlens}[title={What you've built}]
One more word for telling, explaining and describing at B1.
\end{grammarlens}

\subsection*{Your turn}
\begin{itemize}
\raggedright
  \item \emph{Say it:} \emph{pokoléniye}
  \item \emph{Say it:} \emph{pokoléniye}, once more
  \item \emph{Say it:} a sentence of your own with \emph{pokoléniye}
  \item \emph{Recall:} say the Russian for to grow old, then the Russian for a graduation party, then say \emph{pokoléniye} again
\end{itemize}

\begin{tcolorbox}[breakable,colback=teal!4,colframe=teal!35!black,title={Before you move on}]
What does \ru{поколение} mean? (\textquotedblleft{}A generation\textquotedblright{}.) Say it once more.
\end{tcolorbox}
